% Primary and supplementary source basis: :contentReference[oaicite:0]{index=0} :contentReference[oaicite:1]{index=1} :contentReference[oaicite:2]{index=2}

\section{Complex Numbers}

\begin{conceptbox}
    A \textbf{complex number} is written in rectangular form as
    \[
        z = x + iy,
    \]
    where $x,y \in \mathbb{R}$ and $i^2 = -1$.

    The quantities
    \[
        \Re(z)=x, \qquad \Im(z)=y
    \]
    are called the \textbf{real part} and \textbf{imaginary part} of $z$, respectively.

    The \textbf{complex conjugate} of $z$ is
    \[
        z^* = x - iy,
    \]
    and the \textbf{modulus} of $z$ is
    \[
        |z| = \sqrt{x^2+y^2}.
    \]
\end{conceptbox}

Complex numbers extend the real number system in a natural way. In earlier engineering mathematics, variables were often real-valued. Now the independent or dependent variable may be complex. This is important because many engineering tools, especially the Laplace transform and the $Z$-transform, use complex variables.

A useful viewpoint is geometric: the number
\[
    z=x+iy
\]
is represented by the point $(x,y)$ in the complex plane. The horizontal axis is the real axis, and the vertical axis is the imaginary axis.

\begin{conceptbox}
    For
    \[
        z_1=x_1+iy_1, \qquad z_2=x_2+iy_2,
    \]
    the basic operations are
    \[
        z_1+z_2=(x_1+x_2)+i(y_1+y_2),
    \]
    \[
        z_1-z_2=(x_1-x_2)+i(y_1-y_2),
    \]
    \[
        z_1z_2=(x_1x_2-y_1y_2)+i(x_1y_2+x_2y_1),
    \]
    and if $z_2\neq 0$,
    \[
        \frac{z_1}{z_2}
        =
        \frac{(x_1+iy_1)(x_2-iy_2)}{x_2^2+y_2^2}.
    \]
    Also,
    \[
        zz^*=|z|^2=x^2+y^2,
        \qquad
        \Re(z)=\frac{z+z^*}{2},
        \qquad
        \Im(z)=\frac{z-z^*}{2i}.
    \]
\end{conceptbox}

\subsection*{Motivation and Intuition}

Why do we need complex numbers in engineering? A sinusoid such as
\[
    A\cos(\omega t+\phi)
\]
can be handled much more efficiently through
\[
    Ae^{i(\omega t+\phi)}.
\]
Instead of separating magnitude and phase repeatedly, complex numbers allow both to be carried at once in a single expression.

\subsection*{Geometric Interpretation}

If $z=x+iy$, then:

\begin{itemize}
    \item $|z|$ is the distance from the origin to the point $z$.
    \item $z^*$ is the reflection of $z$ across the real axis.
    \item Addition corresponds to vector addition in the plane.
    \item Multiplication has a deeper geometric meaning that becomes clearer in polar form.
\end{itemize}

\subsection*{Key Properties}

\begin{conceptbox}
    Important modulus properties:
    \[
        |z_1z_2|=|z_1||z_2|,
        \qquad
        \left|\frac{z_1}{z_2}\right|=\frac{|z_1|}{|z_2|} \quad (z_2\neq 0),
    \]
    \[
        |z_1+z_2|\le |z_1|+|z_2|
    \]
    (the triangle inequality).
\end{conceptbox}

These are frequently used in complex analysis, signal processing, and control.

\begin{examplebox}
    \textbf{Example 1: Simplify a quotient}

    Simplify
    \[
        \frac{3-2i}{-1+i}.
    \]

    Multiply numerator and denominator by the conjugate of the denominator:
    \[
        \frac{3-2i}{-1+i}\cdot \frac{-1-i}{-1-i}
        =
        \frac{(3-2i)(-1-i)}{(-1+i)(-1-i)}.
    \]

    Compute the numerator:
    \[
        (3-2i)(-1-i)= -3-3i+2i+2i^2.
    \]
    Since $i^2=-1$,
    \[
        -3-3i+2i+2(-1) = -5-i.
    \]

    Compute the denominator:
    \[
        (-1+i)(-1-i)=(-1)^2-i^2 = 1-(-1)=2.
    \]

    Hence,
    \[
        \frac{3-2i}{-1+i}=\frac{-5-i}{2}=-\frac{5}{2}-\frac{1}{2}i.
    \]
\end{examplebox}

\begin{examplebox}
    \textbf{Example 2: Use the conjugate to find real and imaginary parts}

    Let
    \[
        z=4-3i.
    \]

    Then
    \[
        z^*=4+3i.
    \]

    Now,
    \[
        \Re(z)=\frac{z+z^*}{2}
        =
        \frac{(4-3i)+(4+3i)}{2}
        =
        \frac{8}{2}=4,
    \]
    and
    \[
        \Im(z)=\frac{z-z^*}{2i}
        =
        \frac{(4-3i)-(4+3i)}{2i}
        =
        \frac{-6i}{2i}=-3.
    \]
\end{examplebox}

\begin{noteBox}
    In AC circuit analysis, impedance is written as a complex number:
    \[
        Z=R+iX,
    \]
    where $R$ is resistance and $X$ is reactance. This compact form allows algebraic handling of phase shifts and amplitudes in one framework.
\end{noteBox}

\subsection*{Common Mistakes / Notes}

\begin{itemize}
    \item Do not confuse $z^2$ with $|z|^2$. In general,
          \[
              z^2=(x+iy)^2=x^2-y^2+i2xy,
          \]
          whereas
          \[
              |z|^2=x^2+y^2.
          \]
    \item The imaginary part of $z=x+iy$ is $y$, not $iy$.
    \item When dividing complex numbers, always rationalize the denominator using the conjugate.
\end{itemize}